\begin{figure}[H]
\centering
\resizebox{\textwidth}{!}{% Unified architecture (single figure), light-dark palette, serpentine flow.
% Dilation graph + MC-dropout detail embedded inline in the backbone boxes;
% only straight arrows + one turn => no overlapping lines.
\definecolor{acBlueD}{HTML}{1F3B57}
\definecolor{acBlue}{HTML}{3B6EA5}
\definecolor{acBlueL}{HTML}{A6C5E3}
\definecolor{acRed}{HTML}{C0392B}
\begin{tikzpicture}[
    font=\footnotesize, >={Stealth[scale=0.95]},
    box/.style={rectangle, rounded corners=4pt, draw=acBlue, thick, fill=acBlueL!25,
                align=center, inner sep=4pt, text=black},
    dk/.style={box, fill=acBlue, draw=acBlueD, text=white},
    big/.style={rectangle, rounded corners=6pt, draw=acBlueD, very thick, fill=acBlueL!12},
    fl/.style={-{Stealth[scale=1.1]}, very thick, draw=acBlueD},
    fb/.style={-{Stealth[scale=1.0]}, very thick, draw=acRed, densely dashed},
    sub/.style={rectangle, rounded corners=2pt, draw=acBlue, fill=acBlueL!45,
                align=center, font=\tiny, inner sep=2pt, minimum height=0.42cm},
    subd/.style={sub, fill=acBlue, text=white, draw=acBlueD},
    mini/.style={-{Stealth[scale=0.7]}, draw=acBlueD, line width=0.5pt},
]
% ===================== ROW A (left -> right) =====================
\node[box, text width=2.2cm, fill=acBlueL!35, minimum height=1.7cm] (inp) at (0,0)
   {\textbf{Sensor Input}\\ $\mathbf{X}\!\in\!\mathbb{R}^{W\times n}$\\ 5 ch, $W{=}20$};

\node[big, minimum width=4.7cm, minimum height=3.7cm] (tcn) at (4.7,0) {};
\node[anchor=north, font=\bfseries\scriptsize, text=acBlueD, text width=4.5cm, align=center]
   at ($(tcn.north)+(0,-0.07)$) {TCN Backbone --- Dilated Causal Convolutions};
\begin{scope}[shift={(3.5,-0.95)},
   nd/.style={circle, draw=acBlue, fill=acBlueL!40, line width=0.3pt, minimum size=0.17cm, inner sep=0},
   ac/.style={circle, draw=acBlueD, fill=acBlueD, line width=0.3pt, minimum size=0.17cm, inner sep=0},
   rf/.style={-, draw=acRed, line width=0.6pt}]
   \def\dx{0.305}\def\dy{0.40}
   \foreach \r/\lab in {0/$\mathbf{x}$,1/$d{=}1$,2/$d{=}2$,3/$d{=}4$}{
     \foreach \i in {1,...,8} \node[nd] (n\r-\i) at (\i*\dx,\r*\dy) {};
     \node[anchor=east, font=\tiny] at (0.18,\r*\dy) {\lab};}
   \foreach \a/\b in {n2-8/n3-8,n2-4/n3-8,n1-8/n2-8,n1-6/n2-8,n1-4/n2-4,n1-2/n2-4,
                      n0-8/n1-8,n0-7/n1-8,n0-6/n1-6,n0-5/n1-6,n0-4/n1-4,n0-3/n1-4,
                      n0-2/n1-2,n0-1/n1-2} \draw[rf] (\a)--(\b);
   \foreach \n in {n0-1,n0-2,n0-3,n0-4,n0-5,n0-6,n0-7,n0-8,n1-2,n1-4,n1-6,n1-8,n2-4,n2-8,n3-8}
      \node[ac] at (\n) {};
   \node[anchor=west, font=\tiny\bfseries, text=acRed] at ($(n3-8.east)+(0.03,0)$) {$\hat{y}_t$};
\end{scope}
\node[anchor=south, font=\tiny, text width=4.4cm, align=center] at ($(tcn.south)+(0,0.05)$)
   {$k{=}2,\ d_l{=}2^{l-1}{=}\{1,2,4\},\ 64$ ch; $\mathrm{RF}{=}1{+}(k{-}1)\!\sum_l d_l{=}8$};

\node[big, minimum width=3.7cm, minimum height=3.7cm] (mc) at (9.7,0) {};
\node[anchor=north, font=\bfseries\scriptsize, text=acBlueD] at ($(mc.north)+(0,-0.07)$) {MC-Dropout ($M$ passes)};
\begin{scope}[shift={(8.5,-0.55)}]
  \foreach \k/\op in {0/0.5,1/0.7,2/1.0}
     \draw[rounded corners=2pt, draw=acBlue, fill=acBlueL!35, fill opacity=\op]
        (\k*0.18,\k*0.18) rectangle ++(1.25,0.95);
  \node[font=\tiny,align=center] at (0.95,-0.27){active\\dropout};
\end{scope}
\node[font=\scriptsize, text=acBlue] at (10.55,0.0) {$\Rightarrow$};
\begin{scope}[shift={(10.75,-0.4)}]
  \draw[draw=acBlue,fill=acBlueL!35,line width=0.5pt]
    plot[domain=0:1.0,samples=24](\x,{0.42*exp(-((\x-0.5)^2)/0.045)})--(1.0,0)--(0,0)--cycle;
\end{scope}
\node[anchor=south, font=\tiny, text width=3.5cm, align=center] at ($(mc.south)+(0,0.05)$)
   {$\hat{\mathbf{y}}_t{=}\tfrac1M\!\sum_m\!\hat{\mathbf{y}}_t^{(m)},\ \boldsymbol{\sigma}_t{=}\mathrm{std}_m$; $O(M^{-1/2})$ (Thm.~\ref{thm:mc_convergence})};

\node[box, text width=2.5cm, minimum height=1.9cm] (out) at (13.7,0)
   {\textbf{Per-sensor outputs}\\ mean $\hat{\mathbf{y}}_t$ \&\\ epistemic $\boldsymbol{\sigma}_t$};

% ===================== ROW B (right -> left) =====================
% --- AUAA (broken down: multi-head attention + uncertainty gate) ---
\node[big, minimum width=4.1cm, minimum height=3.35cm] (auaa) at (13.7,-5.25) {};
\node[anchor=north, font=\bfseries\tiny, text=acBlueD, text width=3.9cm, align=center]
   at ($(auaa.north)+(0,-0.06)$) {AUAA --- uncertainty-gated attention};
\node[sub, text width=2.9cm] (mha) at (13.7,-4.62) {4-head multi-head attention\\ $\mathrm{softmax}(\mathbf{QK}^\top/\sqrt{d_k})$};
\node[subd, text width=2.9cm] (agate) at (13.7,-5.62) {uncertainty gate\\ $g_t{=}1/(1+\boldsymbol{\sigma}_t)$};
\draw[mini] (agate.north) -- (mha.south) node[midway, right=0pt, font=\tiny, text=acRed]{$\odot$};
\node[anchor=south, font=\tiny, text width=3.9cm, align=center] at ($(auaa.south)+(0,0.04)$)
   {$\boldsymbol{\alpha}^{\mathrm{att}}_t=\mathrm{softmax}(\tfrac{\mathbf{QK}^\top}{\sqrt{d_k}})\odot g_t$};
% --- Dynamic Weight Adapter (broken down: encoders + GRU -> alpha_t) ---
\node[big, minimum width=4.1cm, minimum height=3.35cm] (dwa) at (9.0,-5.25) {};
\node[anchor=north, font=\bfseries\tiny, text=acBlueD, text width=3.9cm, align=center]
   at ($(dwa.north)+(0,-0.06)$) {Dynamic Weight Adapter (GRU)};
\node[font=\scriptsize] (det_e) at (7.55,-4.7) {$\mathbf{e}_t$};
\node[sub, minimum width=0.75cm] (enc1) at (8.35,-4.7) {Enc};
\node[font=\scriptsize] (det_u) at (7.55,-5.45) {$\mathbf{u}_t$};
\node[sub, minimum width=0.75cm] (enc2) at (8.35,-5.45) {Enc};
\node[subd, text width=1.0cm] (gru) at (9.6,-5.05) {GRU\\ $r_t,z_t$};
\node[font=\scriptsize\bfseries, text=acBlueD] (alp) at (10.05,-5.9) {$\alpha_t$};
\draw[mini] (det_e)--(enc1); \draw[mini] (det_u)--(enc2);
\draw[mini] (enc1.east)-- (gru.west); \draw[mini] (enc2.east)-- (gru.west);
\draw[mini] (gru.south)|- (alp.west);
\node[anchor=south, font=\tiny, text width=3.9cm, align=center] at ($(dwa.south)+(0,0.04)$)
   {$\alpha_t=\sigma(\mathbf{W}_\alpha[\mathbf{e}_t,\mathbf{u}_t,\mathbf{h}_t]+\mathbf{b}_\alpha)$};
\node[dk, text width=3.0cm, minimum height=2.0cm] (hyb) at (5.0,-5.25)
   {\textbf{Hybrid Score}\\ $s_t{=}\alpha_t\mathbf{e}_t$\\ $+(1{-}\alpha_t)\mathbf{u}_t$};
\node[box, text width=3.2cm, minimum height=2.0cm] (thr) at (0.8,-5.25)
   {\textbf{Threshold}\\ val-$P_{95}$ \emph{or} Chebyshev\\ $\tau{=}\bar{s}{+}k\sigma_s$, $P(\mathrm{FP})\!\le\!\tfrac1{k^2}$};
\node[dk, text width=2.6cm, minimum height=1.0cm] (det) at (0.8,-7.6)
   {\textbf{Detection}\\ flag if $s_t>\tau$};

% ===================== ARROWS (straight + one turn) =====================
\draw[fl] (inp.east) -- (tcn.west);
\draw[fl] (tcn.east) -- (mc.west);
\draw[fl] (mc.east) -- (out.west);
\draw[fb] (out.south) -- (auaa.north)
   node[midway, right=1pt, font=\scriptsize\bfseries, text=acRed, align=left]{uncertainty\\$\to$ attention};
\draw[fl] (auaa.west) -- (dwa.east);
\draw[fl] (dwa.west) -- (hyb.east);
\draw[fl] (hyb.west) -- (thr.east);
\draw[fl] (thr.south) -- (det.north);
\end{tikzpicture}
}
\caption{End-to-end architecture of the proposed uncertainty-aware TCN framework, organised as a four-stage pipeline: \textbf{Data Pipeline} (loading, leakage-controlled split, windowing) $\to$ \textbf{Model Architectures} (Base TCN and Full Enhanced) $\to$ \textbf{Hybrid Anomaly Scoring} $\to$ \textbf{Anomaly Detection}. The key backbone components are broken down inline: the Base TCN's dilated causal convolution stack (kernel $k{=}2$, dilations $d_l{=}\{1,2,4\}$, receptive field $\mathrm{RF}{=}8$; the red path traces the dilated dependencies feeding the output $\hat{y}_t$), MC-dropout inference ($M$ stochastic passes $\Rightarrow$ predictive mean $\hat{\mathbf{y}}_t$ and per-sensor epistemic $\boldsymbol{\sigma}_t$), the AUAA (4-head attention element-wise modulated by the uncertainty gate $g_t{=}1/(1{+}\sigma_t)$), and the Dynamic Weight Adapter (GRU over $\mathbf{e}_t,\mathbf{u}_t$ producing the learned weight $\alpha_t$). Dark boxes mark the framework's contributed components; the dashed red arrow marks the uncertainty-to-attention coupling that distinguishes this work from approaches that apply uncertainty only post hoc; detection via validation-calibrated thresholding terminates the pipeline.}
\label{fig:architecture}
\end{figure}
